
%%% ====================================================================
%%% ====================================================================
%%%
%%%  SECTION 7: CONCLUSIONS
%%%
%%% ====================================================================
%%% ====================================================================

\section{Conclusions}\label{sec:conclusions}

The PCF framework, formalized and verified in Lean~4,
establishes a closed deductive chain from the decidable
arithmetic of~$\Ztwenty$ to $\mathrm{Re}(\rho)=1/2$.
The chain operates through three
levels---decidable (Verified), structural (spectral parameters from
categorical uniqueness), and analytic (functional equation
from Hecke 1920)---converging at the value~$1/2$ from
independent directions.

The main results:
\begin{itemize}
  \item \lean{rh_squeeze} (\cref{ssec:squeeze}):
    $\mathrm{Re}(\rho)=1/2$ within the PCF categorical setting, via two
    independent bounds neither of which alone suffices.
  \item $T^*$ operator (\cref{sec:results}): ${<}1.7\%$ error across twelve
    orders of magnitude (125~zeros, $n = 1$ to $10^{12}$), with $T^*(n)/t_n\to 1$.
    Convergence properties and bias factors formalized.
  \item \lean{spectral_mapping} (\cref{ssec:zeta-cat}):
    the identification of non-trivial zeros with the ternary fragment is
    derived, not assumed---closing the deductive chain.
  \item The only classical input is Hecke's 1920 theorem (unconditional):
    for self-dual characters, $L(\rho,\chi)=0 \implies L(1-\rho,\chi)=0$.
\end{itemize}

Every spectral invariant ($d=3$, $\mu=1/2$, $\sigma=3/2$) emerges from
PCF's own $S_3$ geometry, with no component of~$\zeta(s)$ entering the
derivation.  The ring $\Rpcf = \mathbb{Z}[\golden,\golden^{-1},\tfrac{1}{2}]$
provides $\mathbb{F}_1$-descent data in the sense of Borger, placing the
construction within Manin's program for absolute geometry.

The proof proceeds by squeeze via two convergent bounds: the universal
ternary bound $\mathrm{Re}(\rho)\le\mu_3=1/2$ (from
\lean{spectral_mapping} $+$ \lean{spectral_uniqueness}) and the
companion bound $\mathrm{Re}(1-\rho)\le\mu_3$ (from
\lean{golden_group_exponent_two} $+$ Hecke's functional equation).  The structural correspondence between $S_3$ invariants and the
Euler product---mediated by the embedding of primes $p>5$ into
$\Ztwenty$ (\cref{thm:zeta-forced})---suggests the
framework has scope beyond the Riemann spectrum.
